\documentclass[12pt,a4paper]{article}

\usepackage[brazilian]{babel}

\usepackage{amsmath}
\usepackage{amssymb}
\usepackage{enumitem}
\usepackage{etoolbox}
\usepackage{tocloft}
\usepackage{tikz}

\setlength{\cftsecnumwidth}{3.2cm}
\setlength{\cftsubsecnumwidth}{1.5cm}

\renewcommand{\thesection}{Capítulo \arabic{section} -}
\renewcommand{\thesubsection}{\arabic{section}.\arabic{subsection}}
\pretocmd{\subsection}{\clearpage}{}{}

\title{
    \Huge \textbf{Um Curso de Cálculo Volume 1} \\
    \vspace{0.3cm}
    \Large 6\textordfeminine{} Edição (Guidorizzi) \\
    \vspace{0.5cm}
    \large \textit{Resoluções Completas e Passo a Passo do Capítulo 6 Funções Exponencial e Logarítmica}
}
\author{Rogério Campos}

\begin{document}

\maketitle
\newpage

\tableofcontents
\newpage

\setcounter{section}{5}
\setcounter{subsection}{0}

\section{Funções Exponencial e Logarítmica}

\subsection{Potência com Expoente Real}
\vspace{0.5cm}
\noindent\textbf{Exercício 1.}
\begin{enumerate}
    \renewcommand{\theenumi}{\alph{enumi}}
    \renewcommand{\labelenumi}{\theenumi)}
    \item 
    \item 
    \item 
    \item 
    \item 
    \item 
    \item 
    \item 
    \item 
    \item 
\end{enumerate}

\vspace{0.5cm}
\noindent\textbf{Exercício 2.}
\begin{enumerate}
    \renewcommand{\theenumi}{\alph{enumi}}
    \renewcommand{\labelenumi}{\theenumi)}
    \item 
    \item 
    \item 
    \item 
    \item 
    \item 
    \item 
    \item 
    \item 
    \item 
\end{enumerate}

\subsection{Logaritmo}
\vspace{0.5cm}
\noindent\textbf{Exercício 1.}
\begin{enumerate}
    \renewcommand{\theenumi}{\alph{enumi}}
    \renewcommand{\labelenumi}{\theenumi)}
    \item 
    \item 
    \item 
    \item 
    \item 
    \item 
    \item 
    \item 
\end{enumerate}

\vspace{0.5cm}
\noindent\textbf{Exercício 2.}
\begin{enumerate}
    \renewcommand{\theenumi}{\alph{enumi}}
    \renewcommand{\labelenumi}{\theenumi)}
    \item 
    \item 
    \item 
    \item 
    \item 
    \item 
\end{enumerate}

\vspace{0.5cm}
\noindent\textbf{Exercício 3.}
\begin{enumerate}
    \renewcommand{\theenumi}{\alph{enumi}}
    \renewcommand{\labelenumi}{\theenumi)}
    \item 
    \item 
    \item 
    \item 
    \item 
    \item 
    \item 
    \item 
\end{enumerate}

\vspace{0.5cm}
\noindent\textbf{Exercício 4.}
\begin{enumerate}
    \renewcommand{\theenumi}{\alph{enumi}}
    \renewcommand{\labelenumi}{\theenumi)}
    \item 
    \item 
    \item 
    \item 
    \item 
    \item 
    \item 
\end{enumerate}

\subsection{O Limite $\lim\limits_{x \to +\infty} \left(1 + \frac{1}{x}\right)^x$}
\vspace{0.5cm}
\noindent\textbf{Exercício 1.}
\begin{enumerate}
    \renewcommand{\theenumi}{\alph{enumi}}
    \renewcommand{\labelenumi}{\theenumi)}
    \item 
    \item 
    \item 
    \item 
    \item 
    \item 
    \item 
    \item 
\end{enumerate}

\vspace{0.5cm}
\noindent\textbf{Exercício 2.}

\vspace{0.5cm}
\noindent\textbf{Exercício 3.}
\begin{enumerate}
    \renewcommand{\theenumi}{\alph{enumi}}
    \renewcommand{\labelenumi}{\theenumi)}
    \item 
    \item 
    \item 
    \item 
\end{enumerate}

\end{document}
